% Evaluar mientras funciona — la línea de tiempo de la roya y las variedades.
%
% Diagrama propio en TikZ (sin librerías extra: una línea, cuatro hitos y una
% llave dibujada a mano con dos segmentos). Se tiñe con el color de la línea
% (milcGradeDark = verde lima de Pensamiento computacional).
%
% Muestra el momento en que se hizo la pregunta: antes de que el problema
% llegara, con los cafetales sanos. Y, debajo, la forma de la respuesta: no
% una variedad sino varias.
%
% Las anotaciones de dentro enseñan; la síntesis va en el pie de figura vía
% \guiaDiagrama, no aquí.
\begin{tikzpicture}[font=\sffamily,>=stealth]

  % ── La línea de tiempo ───────────────────────────────────────────────
  \draw[->,milcNegro!55,line width=1pt] (-0.4,0) -- (11.4,0);

  \foreach \x in {0.6,3.6,6.9,9.9}{
    \draw[milcNegro!55,line width=0.8pt] (\x,-0.16) -- (\x,0.16);
  }

  % ── Hitos ────────────────────────────────────────────────────────────
  \node[fill=milcMostaza,draw=milcNegro,line width=0.5pt,rounded corners=3pt,
        inner sep=4pt,align=center,font=\scriptsize\bfseries,text width=2.4cm]
    at (0.6,-1.0) {antes de 1983\\[0.4mm]\mdseries\itshape los cafetales están sanos};

  \node[fill=milcGradeDark,text=white,rounded corners=3pt,inner sep=4pt,
        align=center,font=\scriptsize\bfseries,text width=2.6cm]
    at (3.6,-1.05) {1983\\[0.4mm]\mdseries\itshape llega la roya · sale la variedad Colombia};

  \node[fill=milcGradeDark,text=white,rounded corners=3pt,inner sep=4pt,
        align=center,font=\scriptsize\bfseries,text width=2.4cm]
    at (6.9,-1.0) {2005\\[0.4mm]\mdseries\itshape empieza la siembra de Castillo};

  \node[fill=milcGradeDark,text=white,rounded corners=3pt,inner sep=4pt,
        align=center,font=\scriptsize\bfseries,text width=2.4cm]
    at (9.9,-1.0) {después\\[0.4mm]\mdseries\itshape Tabi y Cenicafé 1};

  % ── La llave: dónde se hizo la pregunta ──────────────────────────────
  \draw[milcGradeDark,line width=1pt]
    (0.6,0.55) -- (0.6,0.85) -- (3.6,0.85) -- (3.6,0.55);
  \draw[milcGradeDark,line width=1pt] (2.1,0.85) -- (2.1,1.15);
  \node[above,align=center,font=\scriptsize\bfseries,milcGradeDark,text width=5.2cm]
    at (2.1,1.15) {la pregunta se hizo aquí:\\cuando todavía funcionaba};

  % ── La forma de la respuesta ─────────────────────────────────────────
  \node[anchor=west,align=left,font=\scriptsize\itshape,milcNegro!80,text width=10.5cm]
    at (-0.4,-2.35)
    {La respuesta no fue \textbf{una} variedad resistente, sino \textbf{varias}, apoyadas en diversidad genética: contra razas nuevas del hongo, una sola carta no alcanza.};

\end{tikzpicture}
